 \section{Convex functions on the naturals}
 \label{sec:convexity}
 
 The following basic results about existence of limits for integer-valued or doubly integer-valued functions with convexity properties will be very useful in Section~\ref{sec:construct}.
 
 \subsection{One variable}
 
 \begin{lemma}
 Let $g \colon \mathbb{N} \to \mathbb{R}$ be a convex and bounded function.
 Then $\lim_n g(n)$ exists.
 \label{lem:onevarconvexbounded}
 \end{lemma}
 \begin{proof}
 If $g(n) \geq g(n-1)$ for all $n$, then $g(n)$ is monotone decreasing, and since it is bounded below, it has a limit.
 Otherwise, there exists some $n_0$ so that $g(n_0) > g(n_0-1)$. Since $g$ is convex, the sequence $h(n) \coloneqq g(n+1)-g(n)$ is monotone increasing.
 But then, $g(n) \geq g(n-1)$ for all $n \geq n_0$.
 Indeed, it is true for $n=n_0$ by assumption, and if $g(n) \geq g(n-1)$ for $n > n_0$, then since $h$ is monotone increasing, $g(n+1)-g(n) \geq g(n)-g(n-1)$, so adding both inequalities, $g(n+1) \geq g(n)$.
 Thus, $g(n)$ for $n \geq n_0$ is monotone increasing, and since it is upper bounded, it has a limit.
 \end{proof}
 
 
 We say that a function $h \colon \mathbb{N}$ is \emph{$r$-midpoint convex}, for some $r \in \mathbb{N}$, if $h(n-r) + h(n+r) \geq 2 h(n)$ for all $n \in \mathbb{N}$ where the inequality makes sense.
 
 \subsection{Two variable}
 
 We say that a function $g \colon \mathbb{N} \times \mathbb{N} \to \mathbb{R}$ is \emph{$(r,s)$-midpoint convex} if there exists $r,s \in \mathbb{N}$ so that, for all $m$, $g(\cdot, m)$ is $r$-midpoint convex, and for all $n$, $g(n, \cdot)$ is $s$-midpoint convex.
 We say that $g$ is \emph{doubly convex} if it is $(1,1)$-midpoint convex.
 
 \begin{lemma}
   If $g \colon \mathbb{N} \times \mathbb{N} \to \mathbb{R}$ is doubly convex and bounded, then $\lim_{n,m} g(n,m) = \lim_{n} \lim_m g(n,m) = \lim_m \lim_n g(n,m)$ exists.
 \end{lemma}
 \begin{proof}
 The proof is similar to the single-variable case.
 If $g$ is doubly decreasing, then the limit exist since the sequence is bounded.
 Otherwise, there exists $n_0,m_0$ so that $g(n_0,m_0) < g(n_0,m_0+1)$. Then, for fixed $n_0$, $(g(n_0,m))_{m \geq m_0}$ is monotone increasing, by a similar argument as in Lemma~\ref{lem:onevarconvexbounded}. Since it is bounded, the iterated limits and the double limit exists and are equal \DMGtodo[inline]{We need a reference for this.}
 \end{proof}
 
 \begin{lemma}
 Given a $r$-midpoint convex function $g \colon \mathbb{N} \to \mathbb{R}$ which is bounded, $\lim_n g(n)$ exists.
 \label{lem:limexistrmidconvex}
 \end{lemma}
 \begin{proof}
 Let $g$ be a $r$-midpoint convex function, i.e., that $g(n) \le (g(n+r) + g(n-r))/2$ for $r$ larger than some
$r_0$. WLOG, assume $r_0$ is a power of $2$. For $0 \le k
< r_0$, define $g_k(n) \coloneqq g(k + r_0 \cdot n)$. Then $g_k$ is convex and bounded, and thus by Lemma~\ref{lem:onevarconvexbounded} has
a well-defined limit as $n \to \infty$. We must show these limits
agree, for different $k$. To see that the limit of $g_{r_0/2}$ is less or equal than the limit of $g_0$,
note that $2g((3+2n)r_0/2) \le g(nr_0) + g((n+3)r_0)$. We can iterate
this construction to show all the limits are equal, for instance using
$2g(5r_0/4) \le g(0) + g(5r_0/2)$.
 \end{proof}
 
 The proof of the analogous lemma for two-variable functions is very similar to that of Lemma~\ref{lem:limexistrmidconvex}.
 
  \begin{lemma}
 Given a $(r,s)$-midpoint convex function $g \colon \mathbb{N} \to \mathbb{R}$ which is bounded, $\lim_{n, m} g(n,m) = \lim_n \lim_m g(n,m) = \lim_m \lim_n g(n,m)$ exist.
 \label{lem:limexistrmidconvextwovar}
 \end{lemma}